\chapter{Architecture and Software Design}
\label{cap:architecture}
The GeoMedia software architecture is represented by a three tier solution, which divide the whole platform into different actors according they operability and functions.

In specific, a three tier model~\cite{article:three-tier} is composed by the following layers:

\begin{itemize}
    \item \textbf{Client layer (Presentation layer)}: provides the user interface and manages user interactions. It is responsible for presenting data to the user and forwarding user requests to the application layer while remaining independent of the underlying business logic and data storage. \\
    In GeoMedia, this is represented by the mobile application.
    
    \item \textbf{Application layer (Business Logic layer)}: implements the core functionalities of the platform and applies the business logic to process user requests. This layer acts as an intermediary between the presentation and data layers, encapsulating the application's logic and coordinating the exchange of data to ensure consistency and proper execution of system operations.\\
    In GeoMedia, this layer is represented by the backend server which processes client requests and communicates with the data layer.

    \item \textbf{Data layer (Persistence layer)}: is responsible for managing, storage and performing data operations such as the creation, retrieval, update and deletion (CRUD). \\
    In GeoMedia, this layer is represented by the Database Management System (DBMS) that provides a complete solution for data handling.
    
\end{itemize}

This solution provide a classic and simple architecture. By exploiting technology functionalities such network communication with \acrshort{tcpip} stack (analysed in subsection~\ref{sub:tcpipstack}), these actors are able to establish a communication among each other and then to exchange information and activate business process.

The~\acrshort{uml} schema in the Figure~\ref{fig:UML-arch}, graphically represents the established architecture.
\begin{figure}[H]
    \centering
    \includegraphics[width=0.45\textwidth]{UML_Architecture.png}
        \caption{Three-tier Software Architecture UML.}
    \label{fig:UML-arch}
\end{figure}


% @startuml
% top to bottom direction

% package "GeoMedia - Software Architecture" {

%     package "Presentation Tier" {
%         [Mobile Application]
%     }

%     package "Application Tier" {
%         [REST API Server]
%     }

%     package "Data Tier" {
%         [Database Management System]
%     }

%     [Mobile Application] -[hidden]down-> [REST API Server]
%     [REST API Server] -[hidden]down-> [Database Management System]

%     [Mobile Application] --> [REST API Server] : HTTPS (TLS)
%     [REST API Server] --> [Database Management System] : SQL Server Protocol (TDS)\nSQL Queries
% }
% @enduml


Regarding GeoMedia, the communication among the first two layers is implemented through the HTTP protocol~\cite{article:http}, which adoption as the underlying communication protocol of the Web makes it a well-established choice for communication between distributed software components. HTTP is a request-response protocol based on the client-server model that create a virtual link between the application layer (the server) that returns a response to the presentation layer (the clients) with the data requested.
The data, referred to as packets, can be manipulated through intermediary network nodes (e.g., proxy servers, NAT, web caches, etc.) while respecting the TCP/IP structure. To ensure the confidentiality and security of the data, HTTP can be secured through HTTPS, which relies on TLS and digital certificates to authenticate the server and establish an encrypted connection. During the TLS handshake, asymmetric cryptography is used to securely establish a symmetric session key, which is subsequently used to encrypt the data exchanged between the client and server~\cite{site:https-handshake}.

Instead, to establish the connection between the other layer, the protocol depends on the data layer, which in the case of GeoMedia the DBMS is implemented by Microsoft SQL Server (MSSQL) that operate through the TCP/IP stack with $\acrshort{tds}_{\glslink{tdsg}{g}}$ (Tabular Data Stream) protocol.


\section{Microservices Comparison}

Another valid alternative to the architecture adopted is the microservices structure. This paradigm sees the system decomposed into a set of small and independent services, each responsible for a specific business capability~\cite{article:microservices-explained}.

Each microservice encapsulates its own business logic and expose its functionalities through interfaces, usually following the RESTful$_{\glslink{restful}{g}}$ architectural style~\cite{article:rest-api}. REST-based~$\acrshort{api}_{\glslink{apig}{g}}$ support distributed system by adopting constraints such as  stateless communication, cache-ability and uniform interface to support scalable and interoperable systems. This design primarily concerns mainly the backend structure: application and data layer; while the client still represents the single point of interaction with the users.

This solution compared to the a monolithic solution provides high degree of modularity and flexibility and others several key advantages, such as:

\begin{itemize}
    \item \textbf{Scalability:} individual services can be scaled independently according to their workload, reducing infrastructure costs and improving resource utilization.
    
    \item \textbf{Maintainability:} the decomposition into smaller services simplifies both code maintenance than testing and debugging, as each service has a limited and well-defined responsibility.
    
    \item \textbf{Technology independence:} each service can be implemented using the programming language, framework or database that best fits its operability, while respecting the compatibility through communication interfaces.
    
    \item \textbf{Fault isolation:} failures occurring in one service are less likely to compromise the entire application, increasing the reliability of the system.
    
    \item \textbf{Continuous deployment:} updates can be released independently for each service, enabling faster development cycles and reducing deployment risks.
\end{itemize}

Despite these advantages, the adoption of a microservices architecture represent an additional complexity which can be avoided for simple projects or applications with low usage. In fact, the challenge related to network communication, together with issues such as data confidentiality, authentication or service orchestration tend to outweigh the architectural benefits~\cite{article:microsvervices-vs-monolithic}. 

The three-tier architecture designed for GeoMedia represents a monolithic solution that is sufficient for the current size and usage of the application, successfully supporting its functional requirements. Nevertheless, the clear separation between the presentation, application and data layers ensures that the software remains extensible and can progressively evolve towards a microservices architecture without requiring a complete redesign.

On the other hand, the adoption of a microservices architecture could be an interesting and justified scenario to support local and distributed databases across different countries. Since GeoMedia focuses on local data, synchronization would only be required for user accounts and collections, while posts could be stored on country-specific servers according to their geographical coordinates metadata.

A concept of the microservice adoption for GeoMedia is shown in the UML in Figure~\ref{fig:UML-arch-micro}.

\begin{figure}[h]
    \centering
    \includegraphics[width=0.65\textwidth]{UML_Arch_MicroServices.png}
        \caption{Microservices Software Architecture UML.}
    \label{fig:UML-arch-micro}
\end{figure}


% @startuml
% top to bottom direction

% package "GeoMedia - Distributed Microservices Architecture" {

%     package "Presentation Layer" {
%         [Mobile Application]
%     }

%     package "Microservices Layer" {
%         [API Gateway]
%         [Authentication Service]
%         [User Service]
%         [Post Service]
%         [Collection Service]
%     }

%     package "Data Layer - Global" {

%         database "User Database"
%     }

%     package "Data Layer - Local EU/US/.." {

%         database "Post Database"
%         database "Collection Database"
%     }

%     [Mobile Application] -[hidden]down-> [API Gateway]
%     [API Gateway] -[hidden]down-> [Authentication Service]
%     [Authentication Service] -[hidden]down-> "User Database"

%     [Mobile Application] --> [API Gateway] : HTTPS (TLS)

%     [API Gateway] --> [Authentication Service]
%     [API Gateway] --> [User Service]
%     [API Gateway] --> [Post Service]
%     [API Gateway] --> [Collection Service]

%     [Authentication Service] --> "User Database" : SQL
%     [User Service] --> "User Database" : SQL

%     [Post Service] --> "Post Database" : SQL
%     [Collection Service] --> "Collection Database" : SQL
% }

% @enduml